%!TEX root = paper.tex

\section{Background and Overview}
\label{sec:overview}

The paper deals extensively with vectors and matrices.
The $\ell_p$-norm of a vector $\xx \in \Re^{d}$ is defined as
\begin{equation*}
\norm{\xx}_p = \left( \sum_{i = 1}^{d} \left|\xx_i\right|^{p} \right)^{1/p}. 
\end{equation*}
For a matrix $\AA$, we will use $n$ and $d$ to denote its number of
rows and columns respectively, and $\aa_i$ to denote the vector
corresponding to the $i\textsuperscript{th}$ row of $\AA$.
Note that $\aa_i$ is a {\em column} vector.
The $\ell_p$-norm of a vector $\xx$ w.r.t. $\AA$ can then be written as
\begin{equation*}
\norm{\AA \xx}_p = \left( \sum_{i = 1}^{n} \left| \aa_i^T \xx \right|^{p} \right)^{1/p}.
\end{equation*}
We will also assume our matrices are full rank, since otherwise we can
either project onto its rank space, or use pseudo-inverses accordingly.

Most of our analyses revolve around multiplicative errors.
Here we follow the approximation notation from~\cite{CohenLMMPS14:arxiv}.
For a  parameter $\alpha \geq 1$, we say two quantities $x$ and $y$ satisfy
$x \approx_{\alpha} y$ if
\begin{equation*}
\frac{1}{\alpha} x \leq y \leq \alpha x.
\end{equation*}
Note that if $x \approx_{\alpha} y$, then for any power $p$ we have
$x^{p} \approx_{\alpha^{|p|}} y^{p}$.

A cruciual definition in $\ell_2$ row sampling and matrix concentration bounds is
statistical leverage scores.
The statistical leverage score of a row $\aa_i$ is
\begin{equation*}
\ttau_{i}\left(\AA\right)
\defeq \aa_i^T \left( \AA^T \AA \right)^{-1} \aa_i
= \norm{\left( \AA^T \AA \right)^{-1/2} \aa_i}_2^2.
\end{equation*}
It can be viewed as the squared norm of the $i\textsuperscript{th}$ row after the
statistical whitening transform (see e.g.. Hyvarinen et al.~\cite{HyvarinenO00}).
Equivalently, it is also often defined as the squared row norm of the matrix of left singular vectors
of $\AA$.
The following facts about statistical leverage scores underpin their role
in $\ell_2$ row sampling.
\begin{fact}
\label{fact:tau}
\begin{enumerate}
\item \label{item:foster} (Foster's theorem~\cite{Foster53}) $\sum_{i = 1}^{n} \ttau_{i} (\AA) \leq d$,
\item \label{item:upper1} $\ttau_{i}(\AA) \leq 1$.
\item \label{item:l2concentration} ($\ell_2$ Matrix Concentration Bound)
There exists an absolute constant $C_s$ such that for any matrix $\AA$ and any
set of sampling values $\pp_i$ satisfying
\begin{equation*}
\pp_i \geq C_s \ttau_{i} \left( \AA \right) \log{d} \epsilon^{-2},
\end{equation*}
if we generate a matrix $\SS$ with $N = \sum \pp_i$ rows, each chosen
independently as the $i$\textsuperscript{th} basis vector,
times $\frac{1}{\sqrt{\pp_i}}$, with probability $\frac{\pp_i}{N}$
then with high probability we will have
$\norm{\SS \AA \xx}_2 \approx_{1 + \epsilon} \norm{\AA \xx}_2$ for all vectors $\xx$.
\end{enumerate}
\end{fact}
Combining part~\ref{item:foster} and part~\ref{item:l2concentration} immediately implies that replacing
$\AA$ with $O(d \log d / \epsilon^2)$ reweighted row samples from $\AA$ can give a $(1+\epsilon)$ approximation.

$\ell_p$ spaces are more complicated and lack many useful properties of $\ell_2$.
The Lewis weight approach can be seen as a particular way to tap into the niceness of $\ell_2$
by defining, for any matrix $\AA$, a corresponding matrix $\BB$ so that $\norm{\AA \xx}_p$
is in some sense related to $\norm{\BB x}_2$.  One conceivable notion of relatedness would be
to simply minimize the maximum distortion between the norms.  This would define $\BB$ based on
the John ellipsoid for the convex body $\norm{\AA \xx}_p \leq 1$; this is essentially the technique
used in \cite{DasguptaDHKM09} and follow-up works.  However, it does not lead to tight bounds, and
the Lewis approach is different (although Section~\ref{sec:convex} reveals a similarity).

The most naive approach would seem to be to simply set $\BB = \AA$.  Here, however, one
sees that $\BB$ cannot properly capture $\AA$, since $\norm{\BB x}_2$ is not invariant under
``change of density.''  For instance, one may ``split'' a row $\aa_i$ in $\AA$ into $k$ pieces,
each equal to $k^{-1/p} \aa_i$, and $\norm{\AA \xx}_p$ will remain unchanged.
$\norm{\AA \xx}_2$, though, will change, and could be arbitrarily distorted by subdividing different
rows different amounts.

Instead, we adapt this naive approach by using a different ``density.''  We effectively assume
that a given row $\aa_i$ really represents $\ww_i$ rows, each equal to $\ww_i^{-1/p} \aa_i$.
$\ww_i$ may not be an integer, but this does not really matter; they could be viewed as
weights in a weighted $\ell_p$ norm rather than a number of copies.  When switching to $\ell_2$,
the original row $i$ of $\AA$ would still correspond to $\ww_i$ rows equal to $\ww_i^{-1/p} \aa_i$, but
this now has the same effect as one row equal to $\ww_i^{1/2 - 1/p} \aa_i$, rather than simply $\aa_i$
itself.  Putting this together, we will define $\BB = \WW^{1/2 - 1/p} \AA$.

This still leaves the question of how to actually choose the weights $\ww_i$ (the
specific change of density).  Intuitively, we want the split up rows, $\ww_i^{-1/p} \aa_i$
to be normalized in some way.  Lewis's change of density gives a simple and natural
notion of this: each of these normalized rows should have leverage score 1
(defined in terms of a $\ww_i$-weighted $\ell_2$ norm), or, more explicitly,
the $i$\textsuperscript{th} row of $\BB$ should end up with leverage score $\ww_i$.  Note that
this is a somewhat circular characterization: $\ww_i$ must match the leverage
scores of $\BB$, but $\BB$ itself depends on $\ww_i$.  The Gram matrix
of $\BB$ is $\AA^T \WW^{1 - 2/p} \AA$.  Writing this all in terms of the
original matrix $\AA$ gives:
\begin{definition}
\label{def:lewis}
For a matrix $\AA$ and norm $p$, the $\ell_p$ Lewis weights $\wwbar$ are the unique
weights such that for each row $i$ we have
\begin{equation*}
 \wwbar_i = \ttau_{i} \left( \WWbar^{1/2 - 1/p} \AA \right).
\end{equation*}
 or equivalently
\begin{equation*}
\aa_i^T \left( \AA^T \WWbar^{1 - 2/p} \AA \right)^{-1} \aa_i = \wwbar^{2/p}_i.
\end{equation*}
\end{definition}
We will make extensive use of the second formulation
since it groups the $\wwbar_i$ values on one side.
It also involves measuring the operator $\AA^T \WWbar^{1 - 2/p} \AA$ against a fixed
vector $\aa_i$, and therefore allows us to incorporate bounds on the operator.
These definitions have analogs in recent speedups of interior point methods
by Lee and Sidford~\cite{LeeS14}; for example, the requirement on the weight function
in the LP algorithm of~\cite{LeeS13a:arxiv}.

Note that for the case where $p = 2$, $\WWbar^{1/2 - 1/p}$ is the identity matrix,
so the Lewis weights are just the leverage scores.  In the case of general $p$,
it is not immediately clear that the Lewis weights must actually exist or be unique
because of the apparently circularity; existence and uniqueness was first established
by Lewis in~\cite{lewis}.  This paper also includes proofs of the existence and
uniqueness of Lewis weights, which follow directly from the arguments used to
prove our algorithms.  The first is in Section~\ref{sec:compute} (Corollary~\ref{cor:itlewis}
and applies to $p < 4$.  The second is in Section~\ref{sec:convex} (Corollary~\ref{cor:conlewis})
and gives a proof of existence for all $p$ and uniquness for $p \geq 2$ (so together, these
proofs give existence and uniqueness for all $p$); the resulting proof from that section
is a restatement of traditional proofs.

In Section~\ref{sec:concentration} (combined with Section~\ref{sec:ellp}),
we will prove the following concentration bound,
which is a variant of Proposition 2 from~\cite{Talagrand90}.
\begin{restatable}[$\ell_1$ Matrix Concentration Bound]{theorem}{lonechernoff}
\label{thm:l1chernoff}
There is an absolute constant $C_s$ such that given a matrix $\AA$
with $\ell_1$ Lewis weights $\wwbar$, for any set of sampling values $\pp_i$,
$\sum_i \pp_i = N$,
\begin{equation*}
\pp_i \geq C_s \wwbar_i \log(N) \epsilon^{-2},
\end{equation*}
if we generate a matrix $\SS$ with $N$ rows, each chosen
independently as the $i$\textsuperscript{th} standard basis vector,
times $\frac{1}{\pp_i}$, with probability $\frac{\pp_i}{N}$,
then with high probability we have
\begin{equation*}
\norm{\SS \AA \xx}_1
\approx_{1 + \epsilon} \norm{\AA \xx}_1
\end{equation*}
for all vectors $\xx \in \Re^{d}$.  In particular, with constant factor
approximations to the Lewis weights, $O(d \log(d/\epsilon) \epsilon^{-2})$
row samples suffices.
\end{restatable}

The bound on sample count also follows from Fact~\ref{fact:tau}.
As a result, it remains to compute approximate $\ell_p$ Lewis weights.
We present a novel, but extremely simple, iterative scheme that can do this for all $p < 4$.
We repeatedly perform
\begin{equation*}
\ww'_i \leftarrow \left( \aa_i^T \left( \AA^T \WW^{1 - 2/p} \AA \right)^{-1} \aa_i \right)^{p/2}
\end{equation*}
using a possibly approximate algorithm for computing statistical leverage scores.
This routine resembles the use of $\ell_2$-regression to solve $\ell_1$
regression through repeated reweighting~\cite{ChinMMP13}.
Its convergence also mimics the reduction in row counts in the iterative $\ell_p$
row sampling algorithm by Li et al.~\cite{LiMP13}.
We will prove the following Lemma regarding this procedure in Section~\ref{sec:compute}
\begin{lemma}
\label{lem:computeLewis}
For any fixed $p < 4$, given a routine \textsc{ApproxLeverageScores} for computing
$\beta$-approximate statistical leverage scores of rows of matrices of the form $\WW \AA$ for
$\beta = n^{\Omega(\theta)}$, we can compute a $n^{\theta}$ approximation to
$\ell_p$-Lewis weights for $\AA$ with $O(\log(\theta))$ calls to \textsc{ApproxLeverageScores}. 
\end{lemma}

Combining this Lemma with Theorem~\ref{thm:l1chernoff} or \cite{Talagrand90}, using $\theta$ proportional
to $\frac{1}{\log n}$, gives an algorithm satisfying the conditions of Theorem~\ref{thm:main}.
We may use standard techniques, as described, for instance, in \cite{CohenLMMPS14:arxiv}, to
get an input sparsity time algorithm as claimed in the introduction.
The idea is to run in two phases, first using a constant $\theta$, then resparsifying with $\theta$
proportional to $\frac{1}{\log n}$ (getting a $d^{\theta}$ dependence, rather than just $n^\theta$,
turns out to be automatic, since if $n > d^4$, the $\nnz$ term will dominate anyway).
The same process can be used for any $p \leq 4$, with sparsifier
sizes as stated in Section~\ref{sec:ellp}; however, due to the size of the intermediate sparsifier,
the polynomial dependence will be $d^{\max(\omega, p/2+1) + \theta}$.

It is worth noting that the technique in Section~\ref{sec:sparseconvex}, can also be be combined with
these methods (with Lemma~\ref{lem:computeLewis} replacing the convex optimization technique of
Theorem~\ref{thm:polyslow} from Section~\ref{sec:convex}) to give input sparsity time algorithms.
The arguments in Section~\ref{sec:sparseconvex} allow better parameters in the ranges $p \leq 2$
and $p \leq 4$ (this is discussed in more detail there), and this algorithm can support a
slightly better tradeoff between the input sparsity cost and the polynomial term.
